\section{Теоретические основы}

\subsection{Понятие сложности алгоритмов}

Временная сложность алгоритма определяет, как быстро растёт время выполнения при увеличении размера входных данных \cite{knuth1998art}. Для оценки сложности используется асимптотическая нотация «О-большое».

\subsection{Пузырьковая сортировка (Bubble Sort)}

Пузырьковая сортировка является одним из простейших алгоритмов сортировки. Принцип работы заключается в последовательном сравнении соседних элементов и их обмене, если они расположены в неправильном порядке. Каждый проход алгоритма «всплывает» наибольший элемент в конец массива.

Количество сравнений в худшем случае для массива из \texttt{n} элементов описывается формулой:

\begin{equation}
\label{eq:bubble}
T_{\text{пузырьк}}(n) = \frac{n(n-1)}{2} = O(n^2)
\end{equation}

Пространственная сложность пузырьковой сортировки составляет \texttt{O(1)}, так как алгоритм сортирует массив «на месте» без использования дополнительной памяти.

\subsection{Быстрая сортировка (Quick Sort)}

Быстрая сортировка, разработанная Ч. Хоаром в 1960 году \cite{hoare1962quicksort}, использует стратегию «разделяй и властвуй». Алгоритм выбирает опорный элемент \texttt{pivot}, разделяет массив на две части (элементы меньше опорного и элементы больше опорного), после чего рекурсивно сортирует каждую часть.

Среднее количество сравнений описывается рекуррентным соотношением:

\begin{equation}
\label{eq:quick}
T_{\text{быстр}}(n) = 2T_{\text{быстр}}\left(\frac{n}{2}\right) + O(n) = O(n \log n)
\end{equation}

В худшем случае (когда опорный элемент \texttt{pivot} всегда минимальный или максимальный) сложность составляет \texttt{O(n\textsuperscript{2})}, однако на практике такой случай встречается крайне редко.

Пространственная сложность быстрой сортировки в среднем составляет \texttt{O(log n)} из-за рекурсивных вызовов:

\begin{equation}
\label{eq:space}
S_{\text{quick}}(n) = O(\log n)
\end{equation}

\subsection{Сравнительный анализ сложности}

В таблице~\ref{tab:complexity} представлено сравнение сложности алгоритмов.

\begin{table}[h]
\centering
\caption{Сравнение сложности алгоритмов сортировки}
\label{tab:complexity}
\begin{tabular}{|c|c|c|c|}
\hline
\textbf{Алгоритм} & \textbf{Лучший случай} & \textbf{Средний случай} & \textbf{Худший случай} \\
\hline
Пузырьковая & \texttt{O(n)} & \texttt{O(n\textsuperscript{2})} & \texttt{O(n\textsuperscript{2})} \\
\hline
Быстрая & \texttt{O(n log n)} & \texttt{O(n log n)} & \texttt{O(n\textsuperscript{2})} \\
\hline
\end{tabular}
\end{table}

\subsection{Схема алгоритма}

На рис.~\ref{fig:scheme} представлена блок-схема быстрой сортировки.

\begin{figure}[h]
\centering
\includegraphics[width=0.7\textwidth]{scheme.png}
\caption{Блок-схема алгоритма быстрой сортировки}
\label{fig:scheme}
\end{figure}